\section{Register Flag dan Pengaruhnya pada Keputusan}
\label{sec:register-flag}

Register \eng{flag} tidak menyimpan angka yang dihitung, melainkan
keterangan keadaan hasil operasi terakhir. Setiap keterangan
menempati satu bit dan disebut \eng{flag}. Prosesor 8086 memiliki
register 16 bit bernama \reg{FLAGS}; sembilan bit di antaranya
dipakai sebagai penanda.

Contohnya, instruksi \lstasm|ADD AX, BX| menyimpan hasil di \reg{AX},
tetapi prosesor juga mencatat apakah terjadi simpanan keluar, apakah
hasilnya nol, apakah bit tandanya menyala, dan sebagainya. Catatan
itu masuk ke \reg{FLAGS}. Instruksi berikutnya, misalnya lompatan
bersyarat, membaca catatan tersebut untuk mengambil keputusan.

Bagian ini diringkas dari \cite{gunawan2019}, Bagian 2, hlm.~15--18,
dan \cite{sto2001}, BAB IV, hlm.~17--20. Kaitan \eng{flag} dengan
lompatan bersyarat diperkuat oleh \cite{sto2001}, BAB XIII,
hlm.~68--72.

\subsection{Sembilan Flag pada 8086}

Sembilan \eng{flag} tersebut terbagi menjadi dua golongan. Enam
\eng{flag} status ditentukan oleh hasil operasi, sedangkan tiga
\eng{flag} kendali mengatur cara prosesor bekerja. Arti setiap
\eng{flag} dirangkum pada Tabel~\ref{tab:flag-ringkas}.

\begin{table}[H]
    \centering
    \caption{Sembilan \eng{flag} pada 8086 dan kapan bernilai 1.}
    \label{tab:flag-ringkas}
    \small
    \begin{tabular}{@{}cp{4.0cm}p{2.2cm}p{5.4cm}@{}}
        \toprule
        \textbf{Flag} & \textbf{Nama} & \textbf{Kelompok} &
        \textbf{Bernilai 1 bila \dots} \\
        \midrule
        \flag{CF} & \eng{Carry Flag} & Status &
        Operasi tak bertanda menghasilkan simpanan keluar dari bit
        tertinggi \\
        \addlinespace
        \flag{PF} & \eng{Parity Flag} & Status &
        Jumlah bit 1 pada byte rendah hasil operasi genap \\
        \addlinespace
        \flag{AF} & \eng{Auxiliary Flag} & Status &
        Terjadi simpanan keluar dari empat bit rendah \\
        \addlinespace
        \flag{ZF} & \eng{Zero Flag} & Status &
        Hasil operasi bernilai nol \\
        \addlinespace
        \flag{SF} & \eng{Sign Flag} & Status &
        Bit tanda hasil operasi bernilai 1 (hasil negatif) \\
        \addlinespace
        \flag{OF} & \eng{Overflow Flag} & Status &
        Hasil operasi bertanda berada di luar jangkauan \\
        \addlinespace
        \flag{TF} & \eng{Trap Flag} & Kendali &
        Prosesor diminta berhenti setelah setiap instruksi untuk
        \eng{debugging} \\
        \addlinespace
        \flag{IF} & \eng{Interrupt Flag} & Kendali &
        Interupsi yang dapat ditutup diizinkan \\
        \addlinespace
        \flag{DF} & \eng{Direction Flag} & Kendali &
        Operasi string bergerak menurun, dari alamat besar ke kecil \\
        \bottomrule
    \end{tabular}
\end{table}

\eng{Flag} status dibaca oleh program untuk mengetahui keadaan hasil
operasi. \eng{Flag} kendali diatur oleh programmer atau oleh
\eng{debugger} untuk mengubah cara prosesor bekerja. Pada bab ini,
perhatian utama diberikan pada enam \eng{flag} status karena
\eng{flag} tersebut yang paling sering dipakai untuk mengambil
keputusan.

\subsection{Contoh Perubahan Flag pada Penjumlahan}

Cara paling cepat memahami \eng{flag} adalah mengamati beberapa
kasus penjumlahan yang berbeda. Tabel~\ref{tab:contoh-flag-add}
memuat tiga kasus pada register 16 bit.

\begin{table}[H]
    \centering
    \caption{Perubahan \eng{flag} pada tiga kasus penjumlahan 16 bit.
    Tanda ``--'' berarti \eng{flag} bernilai 0.}
    \label{tab:contoh-flag-add}
    \small
    \begin{tabular}{@{}ccccccp{4.6cm}@{}}
        \toprule
        \textbf{Operand 1} & \textbf{Operand 2} & \textbf{Hasil} &
        \flag{CF} & \flag{OF} & \flag{SF} & \textbf{Alasan singkat} \\
        \midrule
        \hex{7FFF} & \hex{0001} & \hex{8000} & -- & 1 & 1 &
        Dua bilangan positif menghasilkan bilangan negatif \\
        \addlinespace
        \hex{8000} & \hex{8000} & \hex{0000} & 1 & 1 & -- &
        Simpanan keluar tak bertanda; dua bilangan negatif
        menghasilkan nol \\
        \addlinespace
        \hex{FFFF} & \hex{0001} & \hex{0000} & 1 & -- & -- &
        Simpanan keluar terjadi, tetapi hasil bertanda tepat nol \\
        \bottomrule
    \end{tabular}
\end{table}

Penjelasan ketiga kasus tersebut adalah sebagai berikut.

\begin{enumerate}
    \item \textbf{\hex{7FFF} $+$ \hex{0001} = \hex{8000}.}
          Sebagai bilangan tak bertanda, hasilnya masih muat dalam
          16 bit sehingga \flag{CF} bernilai 0. Sebagai bilangan
          bertanda, \hex{7FFF} adalah $+32767$ dan hasil \hex{8000}
          ditafsirkan sebagai $-32768$. Dua operand positif
          menghasilkan hasil negatif, sehingga \flag{OF} bernilai 1.
          Bit tanda hasil menyala, sehingga \flag{SF} juga bernilai 1.

    \item \textbf{\hex{8000} $+$ \hex{8000} = \hex{10000}.}
          Sebagai bilangan tak bertanda, hasilnya melebihi 16 bit
          sehingga \flag{CF} bernilai 1. Sisa 16 bit hasilnya adalah
          \hex{0000}, sehingga \flag{ZF} bernilai 1 dan \flag{SF}
          bernilai 0. Sebagai bilangan bertanda, dua bilangan negatif
          menghasilkan nol, padahal jumlah sebenarnya di luar
          jangkauan, sehingga \flag{OF} bernilai 1.

    \item \textbf{\hex{FFFF} $+$ \hex{0001} = \hex{10000}.}
          Simpanan keluar terjadi sehingga \flag{CF} bernilai 1.
          Sebagai bilangan bertanda, \hex{FFFF} adalah $-1$ dan
          \hex{0001} adalah $+1$; hasilnya nol dan masih berada di
          dalam jangkauan, sehingga \flag{OF} bernilai 0 dan
          \flag{ZF} bernilai 1.
\end{enumerate}

Kasus pertama dan kedua menunjukkan bahwa \flag{CF} dan \flag{OF}
menjawab pertanyaan yang berbeda. \flag{CF} berlaku untuk bilangan
tidak bertanda, sedangkan \flag{OF} berlaku untuk bilangan
bertanda. Keduanya tidak boleh dipertukarkan.

\subsection{Pengaruh Flag terhadap Keputusan Program}

\eng{Flag} baru bermakna ketika dipakai untuk mengambil keputusan.
Program dapat memeriksa \eng{flag} setelah operasi aritmatika atau
setelah instruksi \asm{CMP}, lalu memilih jalur eksekusi berdasarkan
keadaan \eng{flag} tersebut. Pasangan \eng{flag} dan lompatan yang
umum dipakai ditunjukkan pada Tabel~\ref{tab:flag-lompat}.

\begin{table}[H]
    \centering
    \caption{Pasangan \eng{flag} dan lompatan bersyarat yang
    memakainya.}
    \label{tab:flag-lompat}
    \small
    \begin{tabular}{@{}cp{4.0cm}p{5.6cm}@{}}
        \toprule
        \textbf{Flag} & \textbf{Lompatan} & \textbf{Makna} \\
        \midrule
        \flag{ZF} & \asm{JE}, \asm{JZ} & Lompat bila hasil nol atau
        kedua nilai sama \\
        \addlinespace
        \flag{ZF} & \asm{JNE}, \asm{JNZ} & Lompat bila hasil tidak
        nol atau kedua nilai berbeda \\
        \addlinespace
        \flag{CF} & \asm{JC}, \asm{JB} & Lompat bila ada simpanan
        keluar atau nilai pertama lebih kecil secara tak bertanda \\
        \addlinespace
        \flag{CF} & \asm{JNC}, \asm{JAE} & Lompat bila tidak ada
        simpanan keluar \\
        \addlinespace
        \flag{SF} & \asm{JS}, \asm{JNS} & Lompat bila bit tanda
        menyala atau tidak \\
        \addlinespace
        \flag{OF} & \asm{JO}, \asm{JNO} & Lompat bila terjadi
        \eng{overflow} bertanda atau tidak \\
        \bottomrule
    \end{tabular}
\end{table}

Rincian cara kerja lompatan bersyarat dan pola perbandingan
\asm{CMP} dibahas pada bab tentang percabangan dan perulangan. Pada
bab ini, yang perlu dipahami adalah bahwa \eng{flag} menyimpan
keadaan hasil operasi dan bahwa keadaan itu dapat dibaca untuk
menentukan langkah berikutnya.

\begin{peringatanbox}
\vspace{1.1em}
Kesalahan yang sering terjadi adalah menukar \flag{CF} dengan
\flag{OF}. \flag{CF} dipakai untuk operasi bilangan tidak bertanda,
sedangkan \flag{OF} dipakai untuk operasi bilangan bertanda. Memakai
\asm{JC} untuk memeriksa hasil bilangan bertanda dapat memberi
kesimpulan yang salah.
\end{peringatanbox}

\begin{infobox}
\vspace{1.1em}
Instruksi pemindahan data seperti \asm{MOV} tidak memperbarui
\eng{flag} hasil operasi aritmatika. Karena itu, \eng{flag} hasil
\lstasm|ADD AX, BX| masih dapat diperiksa walaupun sesudahnya
dijalankan \lstasm|MOV hasil, AX|.
\end{infobox}
